\documentclass{article}
\usepackage{amsmath,hyperref}
\title{A dock changes while an older report is still owed}
\author{}
\date{}
\begin{document}
\maketitle
\paragraph{Operating envelope and scope.}
This is a development family for I1 and L1--L3 in the Cave Circuit 01
simulation data. It tests a joint digital service floor after a
contact loss and an optional second loss of a prepared recipient
before settlement. The proof does not assert physical egress,
unrestricted traffic, an empirical packet rate, natural interoperability,
or independent selection of an acceptance mission.

The older report has fixed, noncompressed size \(a>0\); a second owed
inspection report has size \(b>0\). After the first fault the carrier Q
already has a qualified archive copy and is the only admissible
sender of both reports in the repair window. Both reports must
be credited by \(7/8\) s, both are first eligible in this window,
and Q's budget for the one-second trailing window is \(c\ge b\) bytes.
The active endpoints and their other budgets are fixed.
Separate sender accounting is sufficient for the grants in the
checked cases; any additional traffic on these budgets is excluded.
The positive family requires \(a+b\le c\), delivery of the unique
final grant, and a matching transfer receipt at S if R left its
dock. It permits one stale preparation notice to arrive after the
final grant. Once the final grant is sent, the corresponding
dock and assignment must remain valid through consumption.
These premises restrict the environment and are not sampled
empirical failure assumptions.

\paragraph{Common source chronology.}
A compiled selected
\href{https://github.com/osrf/subt/tree/083c6a8acb0c25aa1a35e34517e628cd3b52ae77}{SubT broker and receive callback}
delivers P's archival copy to Q during staging and validates a
positive P-to-R reference contact at the old site. At the start of
the experiment R is available instead at the prepared alternate site.
P's attempt at time zero actually consumes a transmit allowance
and does not reach R; Q is already positioned beyond that
direct cut. The time and movement of the earlier relocations
are not evaluated.

At \(1/8\), R is still in the prepared dock and sends a historical
preparation receipt. Q delivers the independent inspection packet
at \(1/4\), while this repair is unresolved.
On the second-fault branch, R's position leaves the receiving
dock. The selected segment is 1/64 m along the y axis, with a
piecewise constant acceleration of magnitude 1 in this fixture.
It returns to rest. The dock rule is a declared local positional
threshold; this is not a stock fault detector. In this branch
R reads its own modeled position and sends a named transfer to
S at \(3/8\). S takes over only after its own native callback has
actually received the transfer. R can still receive radio after
leaving the dock, so reachability alone does not qualify it.
At \(5/8\) the currently eligible terminal sends one new
scoped grant to Q. At \(3/4\) an older receipt may be forwarded
late. At \(7/8\) Q reads its scoped callback history, chooses a
terminal from the final grant only, and sends exactly once.
A digital settlement is credit at the eligible receiver,
not packet presence, source preparation, or a confidence flag.

The small shift and resting RF sites were checked against 249,901
triangles of 14 placed meshes and the selected staging fences,
with an expanded per-body box. An integer separation certificate
gives a conservative lower bound of 0.20198 m in that
selection. The native replay uses the actual graph shortest-path
data and native Cave 01 visibility lookup with a chosen 72 m
range, chosen power and fixed random draw. It is not a
whole-world motion, collision, transport or hardware test.

\paragraph{Admission and exact family bound.}
In a positive history the latest valid grant is unique. When
stationary, R's tested local rule continues to qualify the same
recipient; in the moved history, the actual transfer and the
scoped final callback qualify S. Persistence through
consumption is a separate premise. The policy emits no report
from historical preparation alone; the receiver checks
its own authority before crediting a report. Thus it never
credits a second or wrong older settlement in this family.
Both reports fit in the declared window exactly when
\(a+b\le c\). Every result for the positive family follows by
executing the independent report first and the one lawful older
delivery second, not by combining individually successful runs.
If no current grant arrives, Q keeps the independent delivery
and declines the affected step. At the narrow capacity
\(c=a+b\) there is no room for a second older attempt at a
different recipient. Even where both recipients' physical
routes can receive and the custody law permits attempting them,
neither fixed single recipient works across both branches
without the late observation. This is a scope and information
obstruction in the declared operating class, not a theorem
about all routing platforms or message encodings.

A recovery that spends one fresh outgoing byte from Q to query
after the change needs \(a+b+1\le c\), even when granted the same
correct incoming reply for free. Receiving an unsolicited grant
spends one byte at the serving issuer instead; the switched
case also spends one R-to-S transfer byte. The late historical
duplicate spends another old-source byte and no Q transmit byte.
A direct scheduler with the same unsolicited, typed, time-scoped
grant ties. The charged-query and last-arrived-receipt
comparators isolate dependencies; they are not a substitute
for a full state-of-the-art comparison. In particular,
\href{https://arxiv.org/html/2103.11470}{NeBula} already uses data
mules, return-to-contact and managed network traffic. This
study does not claim that its implementation has to poll,
or that it lacks a suitable application protocol.

\paragraph{Boundary and checked correspondence.}
The negative control permits R's departure after its final grant.
Q's tape and transmit allowance in the stable and late-fault
cases are then identical at consumption, but R rejects the
attempt and the old report remains unresolved. Freshness at the
time of issue is insufficient without a hold through consumption
or another qualified recovery step. The earlier inspection still pays.

The executable projection uses the unchanged admission checker
on actual grant callbacks, Q's shared debit, scope and capacity.
The finite source replay sweeps payload sizes, caps, phase of
recipient, omitted enqueue controls, late stale receipts, and
three policies; results are checked against an independently
stated table. Lean checks the Boolean grant refinement,
impossible fixed destination on the two branches and capacity
boundary for general payload sizes given the stated source
abstraction. It does not verify C++ or field dynamics.
\texttt{REPLAY\_NATIVE.txt} gives the exact checks. No I1 or I2 closure
or independent holdout is claimed.
\end{document}
